\section{Hypotheses}
\label{sec:hypotheses}

This section was written on \draftdate, after the code audit of Section~\ref{sec:implementation}
and after two pilot runs on 20 attacked images each, and before the experiments of
Section~\ref{sec:results}. The pilot runs motivated the hypotheses. They are summarised in
Section~\ref{sec:pilot}, and their images are kept out of the experiments
(Section~\ref{sec:data}). The hypotheses are therefore fixed in advance with respect to the
experiments reported here, not with respect to the pilot. The pilot already showed, on 20 images,
part of what every one of them predicts; each rationale below says how much.
\todo{Freeze this section before the confirmatory runs start, with a time stamp that a reader can
check, for example a registration or a tagged commit in a public repository, and record it here.}

We state three confirmatory hypotheses, H1 to H3, each with a rule that says when we regard it as
supported, refuted or undecided, and three exploratory questions, Q1 to Q3, which have no rule.
Table~\ref{tab:rules} collects the rules. Deviations from this plan are reported in
Section~\ref{sec:results}.

\emph{Conventions.} A \emph{cell} is a pair of a surrogate and a detector; with four released
detectors there are 16 cells, four of them white-box and twelve transfer cells. The rules are
applied on the confirmatory set (Section~\ref{sec:data}), to 8-bit files and at the default
threshold unless a rule names another test. Section~\ref{sec:preprocessing} says which analysis
decides. Measures, bounds and error levels are those of Sections~\ref{sec:metrics}
and~\ref{sec:statistics}. A cell or a part of a hypothesis for which a rule gives no verdict
counts as undecided, and a part with no verdict in any cell is undecided. The margin $\delta=0.02$
is the smallest difference in a success rate that we treat as relevant. It is one fifth of the gap
of about ten points that the published ablation shows between its rows 2 and 5
(Table~\ref{tab:paper-ablation}).

\begin{widetable}
\centering
\caption{Decision rules. A part is supported if the condition of the third column holds in every
cell that gets a verdict, and refuted if the condition of the fourth column holds in at least one.
``Upper'' and ``lower'' are one-sided bounds; the number in brackets is the error level of each.
The last column gives when a cell gets no verdict.}
\label{tab:rules}
\small
\setlength{\tabcolsep}{4pt}
\begin{tabular}{@{}l>{\raggedright\arraybackslash}p{0.25\linewidth}>{\raggedright\arraybackslash}p{0.2\linewidth}>{\raggedright\arraybackslash}p{0.2\linewidth}>{\raggedright\arraybackslash}p{0.2\linewidth}@{}}
\toprule
 & Statistic of a cell & Supported if & Refuted if & No verdict if \\
\midrule
H1(a) & discordance of \cfg{A5} and \cfg{A2} &
exact upper [0.025] below $\delta$ & exact lower [0.025/32] above $\delta$ &
fewer than 250 caught images \\ \addlinespace
H1(b) & success rate of \cfg{X2} minus that of \cfg{X4} &
exact upper [0.025] below $\delta$ & exact lower [0.025/32] above $\delta$ &
fewer than 250 caught images \\ \addlinespace
H2(a) & overshoot share of the attacked images, detector equal to surrogate &
Wilson lower [0.025] above 0.8 & Wilson upper [0.025/4] below 0.5 & -- \\ \addlinespace
H2(b) & recovery share, transfer cells &
Wilson lower [0.025] above 0.5 & Wilson upper [0.025/12] below 0.5 &
fewer than 30 evading images \\ \addlinespace
H3 & retention ratio, pooled over the transfer cells, per channel setting &
bootstrap upper [0.025] below 0.5 & bootstrap lower [0.025] above 0.5 &
fewer than 100 evading images in the reference \\
\bottomrule
\end{tabular}
\end{widetable}

\subsection{H1: the attack reaches the saved image through the control image only}
\label{sec:h1}

\emph{Statement.} (a)~In the released configuration, the latent optimisation does not change the
decisions on the saved image: \cfg{A5} and \cfg{A2} agree. (b)~With a clean control image nothing is
left of the attack: under \cfg{X2} the success rate is at the level of the reconstruction floor
\cfg{X4}.

\emph{Rationale.} The control image reaches the output through a path of three convolutions at
full resolution, and the training objective of the \CVAE cannot tell a decoder that uses the latent
from one that copies its control image (Section~\ref{sec:cvae-structure}). The latent optimisation
computes its loss on the plain decoder, not on the skip decoder that produces the saved image
(Section~\ref{sec:released-pipeline}). Either fact alone could make the optimisation ineffective on
the saved image. Part (a) tests the released configuration as it is. Part (b) tests whether the
route through the latent carries anything at all once the pixel-space route is clean. For part (a)
the pilot gave an indirect indication on one surrogate; part (b) was not examined in the pilot.

\emph{Decision rule.} For part (a), the statistic of a cell is the discordance
$\psi_d(\cfg{A5},\cfg{A2})$. Part (a) is supported if its exact upper bound is below $\delta$ in
every cell, and refuted if its exact lower bound is above $\delta$ in at least one cell. For part
(b), the statistic of a cell is $\ASR_d(\cfg{X2})-\ASR_d(\cfg{X4})$. Part (b) is supported if the
upper bound of this difference is below $\delta$ in every cell, and refuted if its lower bound is
above $\delta$ in at least one cell; a negative difference counts as support. A refutation can
come from either part and any of 16 cells, so its bounds are taken at the error level divided by
32. A cell with fewer than 250 caught images gets no verdict. H1 is supported if both parts are
supported and refuted if either part is refuted. Otherwise it is undecided, and we report which
cells decide.

\emph{Reach of the rule.} With 1000 caught images in a cell, part (a) is supported if at most 11
decisions differ and refuted if at least 36 differ; with 600 caught images the two numbers are 5
and 25. Part (b) is supported, with 1000 caught images, if at most 10 images succeed under \cfg{X2}
and not under \cfg{X4} while none goes the other way, and refuted from 48 such images. The region
in which H1 is supported is narrow on purpose: the claim is that the optimisation changes almost
nothing.

\emph{Comparison with the published ablation.} The rule above concerns the released code. To
relate it to the paper we also compute, with \detS as surrogate, the quantity that the released
driver logs, at 5 iterations, the number in the paper, and at 10, the default. For each of the
fifteen transfer values of Table~\ref{tab:paper-ablation} we give the difference between our value
at 5 iterations and the published one, with the interval for the difference of two independent
proportions on 1000 and 800 images, whose half-width is about 4.5 points. We call a value
reproduced if this interval contains zero. Among fifteen values about one is expected to fall
outside by chance, so only a pattern counts: a row whose three values fall outside on the same
side. We call our result inconsistent with the published gap between rows 2 and 5 if, in the same
quantity at 5 iterations, the upper bound of the paired interval for the difference between
\cfg{A5} and \cfg{A2} is below five points, half of that gap, on each of \detE, \detR and \detD.
This comparison is descriptive and has no part in the verdict on H1.

\emph{Evidence per image.} For a given device the pipeline is deterministic up to numerical
effects, so the direct evidence for part (a) is a comparison of the two outputs image by image. We
report the distribution over images of the pixel difference between the outputs of \cfg{A5} and
\cfg{A2} and of the difference of their scores on every detector, each next to the same quantity
for two repeated runs of \cfg{A5} (Section~\ref{sec:harness}).

\emph{Mechanism.} Three measurements say why the outcome is what it is. They enter no rule.
(i)~At each of the four additions in the skip decoder, the norm of the skip features and the norm
of the decoder features to which they are added, from forward passes on test images. (ii)~We decode
the latent of one image with the control image of another and compare the result with the saved
outputs of the two images. (iii)~The correlation of $\xout-\xin$ with $\xtwo-\xin$ and with
$\xone-\xin$, which says which of the two the output follows.

\emph{Further contrasts.} Reported with intervals, without a rule: \cfg{A5} against \cfg{A1}, for
what the pipeline adds to one PGD call; \cfg{A5} against \cfg{X3}, for what the second PGD call
adds; \cfg{A5} against \cfg{X1}, for what the skip decoder changes; \cfg{X5} against \cfg{A2}, for
whether the optimisation reaches the saved image once its loss is computed on the skip decoder;
\cfg{A3} against \cfg{X6}, for whether it works on its own route; and \cfg{A4a} and \cfg{A4b}
against row 4 of the published table.

\emph{Sensitivity.} To see whether part (a) is specific to the released settings, the contrasts
\cfg{A5} against \cfg{A2} and \cfg{X5} against \cfg{A2} are repeated on the development split with
the surrogates \detE and \detS. One setting is varied at a time: 5, 10 or 50 iterations; a
learning rate of $10^{-2}$, $3\times10^{-2}$ or $10^{-1}$; 2, 5 or 10 denoising steps, with as many
inversion steps. This sweep is descriptive. It cannot say what a \CVAE trained with the published
loss would do, because we use the released weights only.

\subsection{H2: decisions flip, scores do not blend in}
\label{sec:h2}

\emph{Statement.} The success rate at the default threshold overstates how well the outputs of the
released pipeline blend in with real images. (a)~On the surrogate, the attacked images overshoot:
their scores lie above almost all scores of real images. (b)~On a transfer detector, most of the
images that evade the default threshold are flagged by a two-sided test calibrated on real images.

\emph{Rationale.} Neither objective of the pipeline is capped or stops at the decision boundary
(Section~\ref{sec:objectives}). On a transfer detector the perturbation can move a score across the
default threshold without moving it into the range of real images. Part (a) is known in general
\pcite{ozbulak2019overoptimized} and was plain in the pilot; the experiment is a replication at
scale. Part (b) was visible in the pilot on two of three transfer detectors and absent on the
third (Section~\ref{sec:pilot}), so it can fail.

\emph{Decision rule.} For part (a), the statistic of a detector is the overshoot share at
$q=0.05$ among all attacked test images, with that detector as surrogate. Part (a) is supported
for a detector if the lower bound of the Wilson interval of this share is above 0.8, and refuted
if the upper bound is below 0.5; we call two sets separable when at least four of five images are
told apart. For part (b), the statistic of a transfer cell is the \emph{recovery share}: among the
caught images that evade the default threshold after the attack, the share that the two-sided test
at $q=0.05$ flags. Part (b) is supported for a cell if the lower bound of the Wilson interval of
this share is above 0.5, and refuted if the upper bound is below 0.5. Cells with fewer than 30
evading images get no verdict. Each part is supported overall if it is supported for every
detector or cell that gets a verdict, and refuted overall if it is refuted for at least one, with
the error level divided by 4 for part (a) and by 12 for part (b). H2 is supported if both parts
are, and refuted if either is.

\emph{Reach of the rule.} With 1000 attacked images, part (a) is supported for a detector from 825
flagged images and refuted up to 460.

\emph{What is printed next to it.} The default threshold and the calibrated test are different
operating points. Next to every recovery share we therefore print the false-positive rates of both
on the real test images and their detection rates on unmodified generated images. A recovery that
is bought with a higher false-positive rate is visible in this way.

\emph{Constructive check.} If part (a) is caused by the uncapped objectives, an attack that stops
at a target score should not overshoot. We test one capped variant: both PGD calls use the step
size $\epsilon/10$ and at most 40 steps and stop at the first step at which the surrogate score
reaches a cap, and the latent optimisation minimises the amount by which the score falls short of
the cap. The cap is the median score of the real development images. \todo{confirm the cap and
the number of steps on the development split} A call that stops early uses less of its bound, so
the capped variant is compared with \cfg{A5} at its measured distortion as well as at its success
rate. This check is exploratory.

\subsection{H3: compression and resizing remove most of the effect}
\label{sec:h3}

\emph{Statement.} After JPEG compression at quality 75, and after down-scaling by 0.5 and
up-scaling back, less than half of the transfer success of the released pipeline remains, when
success is measured at the same false-positive rate before and after.

\emph{Rationale.} A small pixel-level perturbation is attenuated by compression and resampling
\pcite{dziugaite2016jpg,guo2018countering}. The saved image as a whole is not close to its input in
that sense; in the pilot it was 34~dB away. But if H1 holds, the part of the change that carries
the attack is the pixel-level perturbation of the control image. The attack was not evaluated under
post-processing in the paper (Section~\ref{sec:published-eval}), and the released code contains no
such evaluation. At the default threshold the pilot already showed a drop of this size
(Section~\ref{sec:pilot}); it did not measure it at a matched false-positive rate.

\emph{Statistic.} The same processing can change the false-positive rate of a detector
considerably; in the pilot it reached 0.70 for one detector. A comparison at the default threshold
would mix the survival of the attack with this shift. We therefore use, for each channel setting
$h$, the one-sided test calibrated to $q=0.05$ under that setting. For a transfer cell, let
$\Bset'$ be the generated test images that this test catches before the attack both in the
reference setting and under $h$. The \emph{retention ratio} $\rho(h)$ is the number of images in
$\Bset'$ that evade the test under $h$ after the attack, summed over the twelve transfer cells,
divided by the number that evade it in the reference setting, summed in the same way. The
reference setting is the lossless file in the main analysis and JPEG~95 in the format-equalised
analysis (Section~\ref{sec:preprocessing}). The twelve ratios of the single cells are printed next
to the pooled one.

\emph{Decision rule.} H3 has one verdict per channel setting. For JPEG 75 and for resizing by 0.5,
it is supported if the upper bound of the bootstrap interval of $\rho(h)$ is below 0.5, and
refuted if the lower bound is above 0.5. If fewer than 100 images evade in the reference setting
in total, there is no verdict.

\emph{Descriptive companion.} For every detector and channel setting we report the false-positive
rate at the default threshold and the conditional success rate at the default threshold with its
denominator. These numbers describe a defender who does not recalibrate. They enter no rule.

\subsection{Exploratory questions}
\label{sec:questions}

\emph{Q1: detectors from other groups.} How does the conditional success rate on out-of-family
detectors (Section~\ref{sec:threat}) compare with that on within-family transfer detectors? The
paper itself reports lower rates on a CLIP-based detector, between 19\% and 61\% depending on the
surrogate (Section~\ref{sec:published-eval}), and transfer between CNN-based and CLIP-based
detectors is known to be weak \pcite{derosa2024clip}. We expect the same and treat the comparison
as a description of the detectors we could run. Two cautions apply. A handful of detectors does not
represent a family. And a detector whose pre-processing resamples the $224\times224$ file applies a
channel of its own, so that a lower rate may be an effect of that channel.

\emph{Q2: generator.} Do the contrasts of H1 and the effects of the variants differ between the
generators of GenImage, on which the paper says the detectors were trained, and other generators?
In the pilot the detectors missed more unmodified images of the second kind, which shrinks the set
on which a conditional rate is computed (Section~\ref{sec:pilot}). The statistic is the difference
between the paired differences in the two generator sets, with its interval. The comparison is
confounded: Stable Diffusion~2.1 is both a generator outside GenImage and the model inside the
attack.

\emph{Q3: latent and pixel perturbations under the channel.} If H1 holds, the released pipeline
says nothing about whether a perturbation of the latent survives post-processing better than a
perturbation of the pixels. We examine this with the pixel-space route removed: the control image
is clean and the loss is computed on the skip decoder. We add expectation over transformation
\pcite{athalye2018eot}, with differentiable JPEG compression and resizing, to the latent
optimisation. For comparison we add it to PGD, at the same mean LPIPS distortion. We make no
prediction.

\subsection{Pilot observations}
\label{sec:pilot}

Table~\ref{tab:pilot} summarises the two pilot runs. Both use the released configuration with
surrogate \detE and attack 20 generated images. They differ in three respects: the generator of the
inputs, the number of real images, and the device (an Apple laptop with the MPS back end for run A,
an NVIDIA A100 for run B). The generated images were made at 512 pixels from prompts for ten
classes, and the real images are Imagenette photographs, cropped to their central square; both
were resized to $224\times224$. The table is computed from saved 8-bit files at the default
threshold. The scores quoted below for the stages of the pipeline are those that the pilot harness
logs on floating-point tensors. With these sample sizes the uncertainty is large: 14 successes out
of 14 have a Wilson interval of $[0.78,1.00]$. None of these observations is a result.

\begin{widetable}
\centering
\caption{Pilot runs with the released configuration and surrogate \textsf{E} (marked $^\ast$). Run A:
inputs from Stable Diffusion (SD) 2.1, 20 real images. Run B: inputs from SD 1.5, 100 real images.
Rates are given as counts where the counts were recorded; ``n.r.'': not recorded. The rows
``evading'' count only the images caught before the attack under the same setting. The runs
motivated the hypotheses and are not part of the results.}
\label{tab:pilot}
\small
\begin{tabular*}{\linewidth}{@{\extracolsep{\fill}}lrrrrrrrr@{}}
\toprule
 & \multicolumn{4}{c}{Run A (SD 2.1 inputs)}
 & \multicolumn{4}{c}{Run B (SD 1.5 inputs)} \\
\cmidrule{2-5}\cmidrule{6-9}
Setting and quantity & \textsf{E}$^\ast$ & \detR & \detD & \detS & \textsf{E}$^\ast$ & \detR & \detD & \detS \\
\midrule
\emph{Lossless files} \\
Caught before the attack, of 20        & 14 & 12 & 19 & 12 & 18 & 14 & 20 & 20 \\
Evading after the attack               & 14/14 & 8/12 & 11/19 & 10/12 & 18/18 & 9/14 & 9/20 & 20/20 \\
AUC, real against attacked             & 0.00 & 0.80 & 0.94 & 0.98 & 0.00 & 0.74 & 0.96 & 0.98 \\
Real images flagged                    & 0/20 & 0/20 & 0/20 & 0/20 & 0/100 & 1/100 & 1/100 & 0/100 \\
\addlinespace
\emph{JPEG 75} \\
Caught before the attack, of 20        & 12 & 16 & 20 & 19 & 15 & n.r. & n.r. & n.r. \\
Evading after the attack               & 11/12 & 4/16 & 2/20 & 2/19 & 14/15 & 0.07 & 0.05 & 0.10 \\
Real images flagged                    & 0/20 & 4/20 & 3/20 & 14/20 & 11/100 & n.r. & n.r. & n.r. \\
\addlinespace
\emph{Resizing by 0.5} \\
Caught before the attack, of 20        & 20 & 19 & 20 & 20 & n.r. & n.r. & n.r. & n.r. \\
Evading after the attack               & 9/20 & 3/19 & 2/20 & 5/20 & 0.45 & 0.20 & 0.00 & 0.00 \\
Real images flagged                    & 11/20 & 0/20 & 4/20 & 1/20 & n.r. & n.r. & n.r. & n.r. \\
\bottomrule
\end{tabular*}
\end{widetable}

\emph{For H1: scores of the surrogate.} In run A, the first PGD call raises the mean score of \detE
from $-2.0$ to $+114.4$, before truncation to 8 bits. After encoding, inversion, denoising and
decoding with the plain decoder, and before any optimisation, the mean score is $+1.8$, and 7 of
the 20 images are back on the ``generated'' side: the plain round trip removes most of the effect
of the perturbation. Ten iterations of latent optimisation raise the mean score on the plain
decoder to $+88.6$, so the optimisation is effective on the route on which its loss is computed.
On the saved image the mean score is $+153.2$. A second run computes the loss on the skip decoder.
At its first iteration it logs the score of the skip decoder's output for the latent before
optimisation. Image by image, this score and the score of the saved image of the released
configuration differ by 0.49 on average in absolute value and by at most 1.17, for scores between
120 and 207, with no shift in either direction. This is 0.3\% of the score and an upper bound on
what the optimisation as released does to it; truncation to 8 bits alone moves the same score by
0.28 on average, and no repeated run was made.

\emph{For H1: decisions.} After the first PGD call alone, 8 of 12, 11 of 19 and 11 of 12 caught
images evade \detR, \detD and \detS, which is 30 of 43; for the saved image the counts are 8, 11
and 10, which is 29 of 43. On these 20 images the released pipeline thus evaded the transfer
detectors no more often than its own first stage. With the loss on the skip decoder the latent is
not inert: the mean score of the saved image is $166.7$ instead of $153.2$ in run A and $154.8$
instead of $139.7$ in run B. On lossless files this changed none of the 80 decisions in run A;
after JPEG compression or resizing it changed 5 of 160. In run B, on files, 11 instead of 9 of the
14 caught images evade \detR. The tensor log of the same two runs gives 17 instead of 16 of 20
images classified as real by \detR, so at least one decision differs between the tensor and the
file. The saved images are closer to their inputs (PSNR 34.0~dB) than the outputs of the plain
decoder are (24.1~dB). The comparison without optimisation was read from a log and was not run as
configuration \cfg{A2}, and no score of $\xtwo$ was recorded.

\emph{For H2.} In run A, the scores of \detE on the 20 real images lie between 5.5 and 38.4, and
its scores on the 20 attacked files between 121 and 207. On the transfer detectors, the caught
images that evade the default threshold score below every one of the 20 real images in 7 of 8
cases on \detR and in 9 of 10 on \detS, and in none of 11 on \detD, where all evading images lie
inside the range of the real scores.

\emph{For H3.} At the default threshold in run A, the attack evades in 29 of 43 pairs of a caught
image and a transfer detector on lossless files, and in 8 of 55 after JPEG compression; in run B
the rates per detector fall from 0.64, 0.45 and 1.00 to 0.07, 0.05 and 0.10. In run A, \detS flags
14 of the 20 real images after JPEG compression and \detE flags 11 of 20 after resizing. A
false-positive rate of 0.70 on 20 images may also point to a mismatch between our pre-processing
and that of the detector's training, which is one reason for the check of
Section~\ref{sec:preprocessing}.

\emph{For Q2.} The detectors miss more unmodified images from Stable Diffusion~2.1 than from 1.5:
23 against 8 of 80 pairs of an image and a detector (first row of Table~\ref{tab:pilot}). The two
runs were also made on different devices.
